\begin{table}[t]
\centering
\caption{Feature-group ablation on the Mediterranean dataset (spatial-CV out-of-fold, $K=5$, \SI{10}{\kilo\meter} blocks, $n = 77{,}501$). Hyperparameters are fixed at the Optuna-selected values from Table~\ref{tab:hyperparams}; only the feature subset varies. \textbf{Caveat on hyperparameter leakage}: because the hyperparameters were tuned on the full 33-feature stack and re-used (locked) for every ablation variant, the relative RMSE losses reported per dropped feature group should be read as \emph{upper bounds} on the importance of those groups---a variant re-tuned for its reduced feature set might recover some of the lost performance. We accept this trade-off because re-tuning each variant with the same 100-trial Optuna budget would multiply the compute cost by an order of magnitude and would risk conflating feature-effect with tuning-effect. $\Delta\%$ follows the same sign convention as Table~\ref{tab:stratification}: negative values indicate RMSE reduction (improvement). Sentinel-1 SAR backscatter adds no measurable value in the Mediterranean regime; Sentinel-2 optical and hydrological features are the largest contributors. A terrain-only variant (14 features, no satellite or hydrology) recovers $\sim$68\% of the full-stack RMSE gain and represents a viable operational fall-back for settings without satellite access. The baseline RMSE listed here (\SI{2.485}{\meter}) is a refit run with locked hyperparameters and early stopping; the headline RMSE in Section~\ref{sec:overall_performance} (\SI{2.483}{\meter}) comes from the original Optuna training run with the same hyperparameters but a slightly different early-stopping convergence iteration. The \SI{2}{\milli\meter} difference is below the bootstrap CI width of any individual estimate (Section~\ref{sec:overall_performance}) and does not affect interpretation.}
\label{tab:ablation}
\footnotesize
\begin{tabular}{lrrrr}
\toprule
\textbf{Configuration} & \textbf{n feat} & \textbf{RMSE (m)} & \textbf{MAE (m)} & \textbf{$\Delta$\% vs raw} \\
\midrule
All features (baseline)                                & 33 & 2.485 & 1.099 & $-18.6\%$ \\
$-$ Sentinel-1 SAR                                     & 30 & 2.483 & 1.100 & $-18.7\%$ \\
$-$ 3D shape (openness, SVF)                           & 30 & 2.491 & 1.103 & $-18.4\%$ \\
$-$ Landform (MRVBF, MRRTF, geomorphons)               & 30 & 2.489 & 1.101 & $-18.5\%$ \\
$-$ Hydrology (HAND, TWI, flow, valley, stream)        & 25 & 2.549 & 1.121 & $-16.6\%$ \\
$-$ Sentinel-2 optical                                 & 28 & 2.610 & 1.138 & $-14.5\%$ \\
$-$ All satellite (S1 $+$ S2)                          & 25 & 2.616 & 1.146 & $-14.3\%$ \\
Terrain only (no satellite, no hydrology)              & 14 & 2.669 & 1.167 & $-12.6\%$ \\
\bottomrule
\end{tabular}
\end{table}
